\clearpage
\section{Complete normalization of the original private-reference source}
\label{sec:proposition-source-normalization}
This section closes the missing structural bridge. Let $O$ be the actual
opened translated left graph of Section~\ref{sec:proposition-opening},
and let $Q_1$ be the normalized private closure of
Section~\ref{sec:proposition-final-segment}, with its canonical one-input
term sharing tree. We construct a finite structural calculation
\[
 O\ \equiv_{\mathrm{exact}}\ Q_1.
\]
It preserves the private instance, the complete atom tree, the store and
numeral one. All changes occur in the original translated body's sharing
and box structure. No term, translation clause, machine step or observation
predicate is changed.

\subsection{A finite structural calculation with original certificates}
\termdef{proposition-structural-calculation}{Proposition structural calculation}.
A calculation is a finite proof object whose leaves are boundary-fixing
recursive isomorphisms or genuine original filling witnesses for Idem,
BPullW, BPullC, BPerm or the proved sharing-tree relation. Symmetry and
transitivity compose calculations. They do not add transitive closure to
any template relation. Reflexivity is the identity isomorphism, with a
focus-free source. The focus-freedom lemma follows by induction on the
calculation: each leaf proves its two endpoints focus-free; symmetry swaps
them and transitivity keeps the first and last endpoints.
\lean{Copying/PropositionStructuralCalculation}{20}{36}{Structural.Calculation}
\lean{Copying/PropositionStructuralCalculation}{38}{41}{Structural.Calculation.refl}
\lean{Copying/PropositionStructuralCalculation}{42}{55}{Structural.Calculation.focusFree}

The hole-freedom lemma also uses induction. Isomorphisms preserve all labels.
BPullC, BPullW and BPerm preserve the relevant label predicate through their
original fillers. Idem and sharing-tree leaves have hole-free local endpoints,
so the common recursive filling preserves hole freedom in both directions.
Consequently intermediate hole freedom is derived rather than assumed.
An existing weakening calculation embeds leaf by leaf into this larger
calculus without changing its graph endpoints.
\lean{Copying/PropositionStructuralCalculation}{56}{69}{Structural.Calculation.holeFree_iff}
\lean{Copying/PropositionStructuralCalculation}{262}{271}{Structural.Calculation.ofWeakening}

Plugging a calculation into an addressed site lifts every leaf by the
existing genuine-filling constructor. An unchanged serial prefix or suffix
is retained by an identity filling on that component. Changing both serial
components is achieved by changing the first and then the second. The three
tensor constructions use the same argument, preserving branch order.
\lean{Copying/PropositionStructuralCalculation}{70}{83}{Structural.Calculation.plugAt}
\lean{Copying/PropositionStructuralCalculation}{84}{97}{Structural.Calculation.serialRight}
\lean{Copying/PropositionStructuralCalculation}{98}{111}{Structural.Calculation.serialLeft}
\lean{Copying/PropositionStructuralCalculation}{112}{115}{Structural.Calculation.serial}
\lean{Copying/PropositionStructuralCalculation}{117}{128}{Structural.Calculation.tensorRight}
\lean{Copying/PropositionStructuralCalculation}{130}{141}{Structural.Calculation.tensorLeft}
\lean{Copying/PropositionStructuralCalculation}{143}{148}{Structural.Calculation.tensor}

Boxing lifts every leaf into the same box, fixing its bound prefix and free
interface. Equal-interface transport changes no graph data. A typed output
permutation is applied to both endpoints through an actual permuted filling
context; it is not asserted to be an ordered isomorphism. Parallel-left and
parallel-right lifting retain a specified boundary shuffle. Their composition
therefore lifts both operands through Figure~16's exact paired-subterm
constructor. Application adds its unchanged head and complete distributor.
Abstraction adds its unchanged lambda head, box and outer single-use rows;
the payload variant adds the head and box only.
\lean{Copying/PropositionStructuralCalculation}{149}{162}{Structural.Calculation.boxed}
\lean{Copying/PropositionStructuralCalculation}{163}{177}{Structural.Calculation.reboundary}
\lean{Copying/PropositionStructuralCalculation}{272}{285}{Structural.Calculation.permuted}
\lean{Copying/PropositionStructuralCalculation}{311}{330}{Structural.Calculation.parallelRight}
\lean{Copying/PropositionStructuralCalculation}{332}{354}{Structural.Calculation.parallelLeft}
\lean{Copying/PropositionStructuralCalculation}{355}{361}{Structural.Calculation.pair}
\lean{Copying/PropositionStructuralCalculation}{363}{367}{Structural.Calculation.application}
\lean{Copying/PropositionStructuralCalculation}{286}{298}{Structural.Calculation.abstraction}
\lean{Copying/PropositionStructuralCalculation}{299}{310}{Structural.Calculation.payload}

One-step allowance padding raises each original certificate using that
particular template's proved padding-closure theorem. Iterated padding is
induction on the added allowance. Padding to $N$ iterates $N-d$ times and
transports the resulting equality of depth indices. No semantic padding
law is inferred from machine determinism.
\lean{Copying/PropositionStructuralCalculation}{178}{188}{Structural.Calculation.raise}
\lean{Copying/PropositionStructuralCalculation}{190}{193}{Structural.Calculation.pad}
\lean{Copying/PropositionStructuralCalculation}{195}{205}{Structural.Calculation.padTo}

For a closed calculation, exact-count soundness follows by induction from
the already proved robustness consequences of its five leaf families.
The isomorphism case uses a reflexive filling with an isomorphic endpoint.
Symmetry reverses the two refinements; transitivity composes them, using
the derived hole-freedom invariant for the intermediate graph. Plugging
this result into every original binding-free closing context proves
contextual exact equivalence. Unfolding the two refinements gives equality
of the existence of successful runs at each exact count $n$, including
cases in which neither side terminates.
\lean{Copying/PropositionStructuralCalculation}{206}{237}{Structural.Calculation.exact_equiv}
\lean{Copying/PropositionStructuralCalculation}{238}{245}{Structural.Calculation.contextual_exact}
\lean{Copying/PropositionStructuralCalculation}{246}{261}{Structural.Calculation.termination_count}

\subsection{The inner read body and its actual sharing rows}
Write $W_d$ for a term weakening, $C_d$ for the complete canonical one-use
term tree, $U_d=W_d;C_d$ and $V_d=C_d;C_d$. Here $C_d$ contains both its
weakening and contraction; it is never replaced by an identity wire.
The scoped read term is $!x$, with $x$ at index two beneath $w$ and $f$.
For zero-input term trees $A,B$ and a one-input tree $T$, define
\[
 \mathsf{readParts}(A,B,T)=(A\otimes B)\otimes(\mathsf{deref};T).
\]
Its ordered outputs are $w,f,x$. The factored body chooses $A=B=U_0$,
$T=V_0$; the reduced body chooses $A=B=W_0$, $T=C_0$.
\lean{Copying/PropositionSourceGraphs}{11}{11}{Source.readTerm}
\lean{Copying/PropositionSourceGraphs}{12}{12}{Source.weak}
\lean{Copying/PropositionSourceGraphs}{13}{13}{Source.use}
\lean{Copying/PropositionSourceGraphs}{14}{14}{Source.unusedRow}
\lean{Copying/PropositionSourceGraphs}{15}{15}{Source.usedRow}
\lean{Copying/PropositionSourceGraphs}{17}{20}{Source.readParts}
\lean{Copying/PropositionSourceGraphs}{22}{22}{Source.readFactored}
\lean{Copying/PropositionSourceGraphs}{23}{25}{Source.readReduced}

The first whole-graph certificate identifies the actual eleven-edge
translation with that factored graph. Its inverse maps, every label, and
every ordered incidence are checked against the complete finite inventories.
The unused row has three edges and ranks $(0,0,1)$; the used row has four
edges and ranks $(0,1,0,2)$. Exhaustive labels show that they contain only
term weakenings and contractions. The ranks strictly increase at every
edge adjacency, so every path has distinct edges. These are genuine
zero-input and one-input contraction trees. The sharing-tree template
therefore replaces $U_0$ by a weakening and $V_0$ by $C_0$. Tensor and serial
lifting yield the reduced read body from the unchanged translation.
\lean{Copying/PropositionSourceGeometry}{27}{56}{Source.ReadFactor.iso}
\lean{Copying/PropositionSourceInner}{18}{28}{Source.unused_tree}
\lean{Copying/PropositionSourceInner}{29}{37}{Source.used_tree}
\lean{Copying/PropositionSourceInner}{39}{43}{Source.unusedCalculation}
\lean{Copying/PropositionSourceInner}{44}{46}{Source.usedCalculation}
\lean{Copying/PropositionSourceInner}{48}{55}{Source.readCalculation}

\subsection{Moving only free resources through the two boxes}
\termdef{private-source-free-row}{Private-source free row}.
For a body $G$ with ordered outputs $w,x$, adjoin an unused free output
last, giving $w,x,z$, and apply the free-output swap extended by the identity
on $w$. Call the result $\mathsf{beforeX}(G)$, with outputs $w,z,x$.
Contract its final two free outputs to obtain $\mathsf{insideUse}(G)$.
This is precisely the canonical one-use row on $x$ inside the box.
The permutation pair records and proves the fixed bound output explicitly.
The specialized read body is $\mathsf{insideUse}(\mathsf{readBody}_0)$,
and its lambda payload is the corresponding smaller read closure.
The reduced read body is isomorphic to $\mathsf{beforeX}$ of this specialized
body, with the additional unused output representing $f$.
\lean{Copying/PropositionSourceGraphs}{26}{28}{Source.swapPair}
\lean{Copying/PropositionSourceGraphs}{29}{33}{Source.beforeX}
\lean{Copying/PropositionSourceGraphs}{34}{36}{Source.insideUse}
\lean{Copying/PropositionSourceGraphs}{38}{38}{Source.readUsed}
\lean{Copying/PropositionSourceGraphs}{40}{40}{Source.smallUsed}
\lean{Copying/PropositionSourceGeometry}{57}{86}{Source.ReadUnused.iso}

The auxiliary endpoints retain the literal rule shapes. The reduced inner
translation is its original abstraction constructor. The contraction-pull
endpoint is a lambda head followed by BPullC's literal right graph.
The serial contraction endpoint places the same two-output lambda payload
before a primitive contraction. The permuted-box endpoint is BPerm's
right graph with its unchanged lambda head. The outside-weakening endpoint
replaces that box by BPullW's literal right graph. The intermediate with
outer rows retains Figure~16's entire two-resource single-use distributor.
The normalized inner graph has an outside weakening for $f$, the original
small read closure followed by $C_1$, and both original outer rows.
\lean{Copying/PropositionSourceGraphs}{42}{44}{Source.innerReduced}
\lean{Copying/PropositionSourceGraphs}{45}{49}{Source.pullRight}
\lean{Copying/PropositionSourceGraphs}{50}{51}{Source.contractBox}
\lean{Copying/PropositionSourceGraphs}{53}{55}{Source.permutedBox}
\lean{Copying/PropositionSourceGraphs}{57}{61}{Source.outsideWeak}
\lean{Copying/PropositionSourceGraphs}{62}{62}{Source.innerAfterWeak}
\lean{Copying/PropositionSourceGraphs}{64}{68}{Source.innerNormal}

The generic unused-output calculation first applies the literal BPerm pair
to move the free swap outside the box, then applies BPullW under that same
external swap. The original bound output stays fixed throughout. The
generic contraction calculation applies literal BPullC to
$\mathsf{beforeX}(G)$ under the unchanged lambda head. Both require only
focus freedom of $G$, which is proved for each concrete use.
\lean{Copying/PropositionSourceInner}{56}{71}{Source.removeWeak}
\lean{Copying/PropositionSourceInner}{72}{81}{Source.pullContraction}

For the small read closure, the first finite certificate reassociates
BPullC's literal output contraction into the serial contraction endpoint.
The unused-output calculation then moves the weakening outside. The second
finite certificate identifies the resulting weakening/contraction with the
complete canonical one-use tree outside the small closure. A further
certificate identifies the outside unused $f$ with a tensor of its weakening
and the same small closure. All three certificates preserve the whole
recursive box label. Combining them with the read-body calculation proves
the original inner lambda equivalent to the normalized inner graph, while
retaining both outer translation rows.
\lean{Copying/PropositionSourceGeometry}{87}{116}{Source.SmallPull.iso}
\lean{Copying/PropositionSourceGeometry}{147}{176}{Source.SmallUse.iso}
\lean{Copying/PropositionSourceGeometry}{207}{215}{Source.UnusedOutside.iso}
\lean{Copying/PropositionSourceInner}{82}{93}{Source.smallUseCalculation}
\lean{Copying/PropositionSourceInner}{94}{106}{Source.innerCalculation}

\subsection{The callback application and the outer private closure}
The callback term is literally $f\,()$ in the environment $f,x$; the
sequence term applies the original read lambda to that callback. The
normalized sequence changes only the function subgraph by the preceding
calculation and retains the actual padded callback and application distributor.
\lean{Copying/PropositionSourceGraphs}{69}{69}{Source.callbackTerm}
\lean{Copying/PropositionSourceGraphs}{70}{70}{Source.sequenceTerm}
\lean{Copying/PropositionSourceGraphs}{72}{75}{Source.sequenceNormal}
\lean{Copying/PropositionSourceNormalization}{64}{69}{Source.sequenceCalculation}

Its two complete sharing components are
\[
\begin{aligned}
 F&=(U_1\otimes((C_1\otimes W_1);D_{1,2}));D_{1,2},\\
 X&=(V_1\otimes((W_1\otimes W_1);D_{1,2}));D_{1,2}.
\end{aligned}
\]
Here $D_{1,2}$ denotes the original canonical two-use term tree at the
same allowance. Both $F$ and $X$ have twelve edges. The schematic sequence
places $X$ after the small read closure and $F$ in the callback's function
branch, retains both application heads and unit, then swaps the outputs
into the order $f,x$. A checked whole-graph isomorphism factors the actual
29-edge normalized sequence into this schematic sequence. Each extracted
tree has exhaustive sharing labels and an explicit strictly increasing
edge rank. Thus both are replaced through the proved sharing-tree relation
by $C_1$, retaining the entire operational skeleton.
\lean{Copying/PropositionSourceGraphs}{76}{79}{Source.callbackTree}
\lean{Copying/PropositionSourceGraphs}{80}{83}{Source.privateTree}
\lean{Copying/PropositionSourceGraphs}{85}{90}{Source.sequenceParts}
\lean{Copying/PropositionSourceOuterGeometry}{27}{56}{Source.SequenceFactor.iso}
\lean{Copying/PropositionSourceNormalization}{15}{28}{Source.callback_tree}
\lean{Copying/PropositionSourceNormalization}{29}{42}{Source.private_tree}
\lean{Copying/PropositionSourceNormalization}{53}{55}{Source.callbackCalculation}
\lean{Copying/PropositionSourceNormalization}{56}{58}{Source.privateCalculation}

The reduced schematic sequence is isomorphic to
$\mathsf{insideUse}(\mathsf{callReadBody}_0)$. Hence the generic BPullC,
BPerm and BPullW calculations apply to the large box as well. Their two
finite endpoint certificates retain its complete nested read box. The
large closure has its free $x$ row moved outside while $f$ remains bound.
\lean{Copying/PropositionSourceOuterGeometry}{57}{86}{Source.SequenceUse.iso}
\lean{Copying/PropositionSourceNormalization}{70}{83}{Source.sequenceSharing}
\lean{Copying/PropositionSourceGeometry}{117}{146}{Source.LargePull.iso}
\lean{Copying/PropositionSourceGeometry}{177}{206}{Source.LargeUse.iso}
\lean{Copying/PropositionSourceNormalization}{84}{96}{Source.largeUseCalculation}

The outer abstraction's own single-use row remains after that moved row.
A representation certificate identifies the resulting suffix with $V_2$;
the intermediate body-normal graph is the large closure followed by $V_2$.
The four-edge rank proof certifies $V_2$ as a one-input tree, so another
sharing leaf replaces it by $C_2$. This proves the full original translated
let body equivalent to the large schematic read closure followed by its
canonical one-use tree.
\lean{Copying/PropositionSourceGraphs}{92}{92}{Source.bodyNormal}
\lean{Copying/PropositionSourceOuterGeometry}{87}{116}{Source.AbstractionSuffix.iso}
\lean{Copying/PropositionSourceNormalization}{43}{51}{Source.outer_tree}
\lean{Copying/PropositionSourceNormalization}{59}{63}{Source.outerCalculation}
\lean{Copying/PropositionSourceNormalization}{97}{106}{Source.bodyCalculation}

The opened whole-program certificate identifies $O$ with the raised original
body followed by the unchanged private cell. Raising the body calculation
uses the proved closure of each leaf family. Padding the schematic closure
and its row commutes by the existing constructor isomorphisms. Retaining
the private cell as a serial suffix therefore yields $Q_1$ exactly. This
completes the source-normalization bridge, with no assumption of the desired
equivalence and no replacement of the allocated store.
\lean{Copying/PropositionSourceOuterGeometry}{117}{125}{Source.Opened.iso}
\lean{Copying/PropositionSourceNormalization}{107}{113}{Source.normalized_raise}
\lean{Copying/PropositionSourceNormalization}{114}{122}{Source.openedCalculation}
